% ==========================================================================
%  Section 6 -- Uniform moment bound and tightness.  Owner: W06.
%  Frozen label: sec:tightness.
%
%  Contents: the probe notation and the H^{-2} energy; the uniform trace
%  constant (lem:trace-constant), derived here in the four steps (i)--(iv)
%  of its proof and quoted from nowhere -- this closes ledger GAP R1; the
%  uniform moment (thm:uniform-moment); the two zero-mode devices
%  (lem:zero-mode); compactness of H^{-2}-bounded sets inside the carrier
%  (thm:compact-tail); tightness (thm:tightness); the external extraction
%  theorem (ext:prokhorov).  Nothing here concerns the marginals of mu_N,
%  the identification of a limit, or any measure other than mu_N.
% ==========================================================================
\section{Uniform moment bound and tightness}\label{sec:tightness}

This section supplies the compactness input that Prokhorov's theorem
consumes, and it derives the uniform trace constant rather than quoting it.
The constant is \cref{lem:trace-constant}, the moment it feeds is
\cref{thm:uniform-moment}, the compactness of the sets that moment controls
is \cref{thm:compact-tail}, and Markov's inequality combines the two into
tightness in \cref{thm:tightness}; \cref{lem:zero-mode} separates the two
weight conventions the argument needs. Nothing here concerns the marginals
of $\mu_N$ or a limit.

Fix $L>0$, $m>0$ and $N\in\N$, $a>0$ with $aN=L$ (\cref{hyp:volume});
along the sequence $N=n+1$ and $a=L/N$. Let
$\cos_k(x)=\cos(2\pi\,k\cdot x/L)$ and
$\sin_k(x)=\sin(2\pi\,k\cdot x/L)$ be the cosine and sine tests of mode
$k\in\Z^2$, and extend the test pairing $\varphi(f)$ of a lattice field,
which is real for real $f$, to complex tests by complex linearity. The
\emph{probe} of the field at $k$ is
\begin{equation}\label{eq:probe}
  \widehat\varphi(k):=\varphi(\cos_k)-i\,\varphi(\sin_k)\in\C ,
\end{equation}
that is, the field paired with the complex exponential of mode $k$ in the
normalisation \eqref{eq:fourier-coefficient}. Only
$\abs{\widehat\varphi(k)}^2=\varphi(\cos_k)^2+\varphi(\sin_k)^2$ enters
below. The hat records the mode; it is not the weighted family
$\widehat c$ of \cref{def:carrier-hminus3}, which records the weight. At
$k=0$, $\cos_0\equiv1$ and $\sin_0\equiv0$, so $\widehat\varphi(0)
=\varphi(1)$. The stored coordinate of a carrier point at $k$ is
$\sqrt{w_3(k)}$ times the probe (\cref{eq:lift}), while the $\HL{2}$ norm below
is written in the coefficient family of \cref{rem:two-families} and carries the
$\widetilde w_2$ weight itself, so the two weights are never both applied; it is
the $\widetilde w_2$-weighted sum of squared probe moduli,
\begin{equation}\label{eq:energy}
  \norm{\varphi}^2_{\HL{2}}
  =\sum_{k\in\Z^2}\widetilde w_2(k)\abs{\widehat\varphi(k)}^2
  \in[0,\infty],
\end{equation}
finite on the $\HL{2}$ subspace and $+\infty$ outside it: the norm of
\cref{def:carrier-hs} with the full weight.

The lattice covariance is an input object here, and the only property of it
that this section uses is its sampled covariance form
\cref{eq:sampled-covariance}. For bounded continuous tests $f,g$ it reads
$\E_{\mu_N}[\varphi(f)\varphi(g)]
=a^{2}\ip{f}{A_N^{-1}g}_{\mathrm{counting}}$, in the counting inner product
\eqref{eq:counting-inner-product} and with $A_N=Q$ the operator of
\cref{def:lattice-laplacian}. Its symbol is at least $m^2$
(\cref{prop:operator}), so $A_N\ge m^{2}$, and for every such $f$
\begin{equation}\label{eq:per-test}
  \E_{\mu_N}\bigl[\varphi(f)^2\bigr]
  \le a^{2}m^{-2}\sum_{x\in I_N}f(ax)^2
  \le m^{-2}(aN)^{2}\norm{f}_\infty^{2}
  =m^{-2}L^{2}\norm{f}_\infty^{2},
\end{equation}
the last identity being \eqref{eq:fixed-volume}.

\subsection{The uniform trace constant}

\begin{lemma}[The uniform $\HL{2}$ trace constant]\label{lem:trace-constant}
Let $L>0$, $m>0$ and
\begin{equation}\label{eq:trace-constant}
  C(L,m):=m^{-2}L^{2}+2\,m^{-2}L^{2}\sum_{k\in\Z^2}w_2(k).
\end{equation}
Then $C(L,m)$ is a finite real number with
$m^{-2}L^{2}\le C(L,m)\le 25\,m^{-2}L^{2}$, and for every $N$, $a$ with
$aN=L$,
\begin{equation}\label{eq:trace-bound}
  \sum_{k\in\Z^2}\widetilde w_2(k)\,
  \E_{\mu_N}\bigl[\abs{\widehat\varphi(k)}^{2}\bigr]\le C(L,m).
\end{equation}
The summand of \eqref{eq:trace-bound} at $k=0$ equals $m^{-2}L^{2}$
exactly; each summand with $k\ne0$ is at most $2\,m^{-2}L^{2}w_2(k)$.
\end{lemma}

\begin{proof}
\emph{(i) The per-mode second moment.} For $k\ne0$ the tests $\cos_k$ and
$\sin_k$ are bounded by $1$ in sup norm, so \eqref{eq:per-test} gives
$\E[\varphi(\cos_k)^2]\le m^{-2}L^{2}$ and
$\E[\varphi(\sin_k)^2]\le m^{-2}L^{2}$; by \eqref{eq:probe} and linearity
in each real coordinate,
\[
  \E\bigl[\abs{\widehat\varphi(k)}^{2}\bigr]
  =\E\bigl[\varphi(\cos_k)^2\bigr]+\E\bigl[\varphi(\sin_k)^2\bigr]
  \le 2\,m^{-2}L^{2}=2Vm^{-2}.
\]
The factor $2$ is exactly the pair of real coordinates: the squared modulus
of a complex random variable is the sum of the squares of its real and
imaginary parts, and each part is a real Gaussian image of the field with
variance at most $Vm^{-2}$, because $\lambda_N\ge m^{2}$.

\emph{(ii) The zero mode, exactly.} At $k=0$ the sine test vanishes and the
cosine test is the constant $1$, so
$\E[\abs{\widehat\varphi(0)}^{2}]=\E[\varphi(1)^2]
=a^{2}\ip{1}{A_N^{-1}1}_{\mathrm{counting}}$ by \eqref{eq:sampled-covariance}.
Constants are annihilated by the discrete Laplacian
(\cref{def:lattice-laplacian}), so $A_N1=m^{2}1$ and $A_N^{-1}1=m^{-2}1$,
while the counting inner product of $1$ with itself is the number of sites
$N^{2}$. Hence
\[
  \E\bigl[\abs{\widehat\varphi(0)}^{2}\bigr]
  =a^{2}m^{-2}N^{2}=m^{-2}(aN)^{2}=m^{-2}L^{2}=V/m^{2},
\]
an equality and not an estimate: it is the covariance pairing of the
constant test with itself. Only the cosine coordinate survives at $k=0$,
which is why this term carries no factor $2$.

\emph{(iii) Summability, with a value.} The shell
$\{\norm{k}_\infty=q\}$ carries $8q$ modes, on which $w_2=q^{-4}$
(\cref{ext:summability}), hence
\[
  \sum_{k\in\Z^2}w_2(k)=\sum_{q\ge1}8q\cdot q^{-4}
  =8\sum_{q\ge1}q^{-3}
  \le 8\Bigl(1+\int_1^\infty t^{-3}\,\dd t\Bigr)=12 ,
\]
the summand at $k=0$ contributing nothing. So the series in
\eqref{eq:trace-constant} converges, and only positivity of $m$ and $L$
enters this step.

\emph{(iv) Assembly.} Since $\widetilde w_2(0)=1$ and
$\widetilde w_2(k)=w_2(k)$ for $k\ne0$ (\cref{def:modes}) and probe moments
are nonnegative, splitting at the zero mode and using (ii) on the first
piece and (i) termwise on the second,
\[
  \sum_{k\in\Z^2}\widetilde w_2(k)\E\bigl[\abs{\widehat\varphi(k)}^{2}\bigr]
  =\E\bigl[\abs{\widehat\varphi(0)}^{2}\bigr]
   +\sum_{k\neq0}w_2(k)\E\bigl[\abs{\widehat\varphi(k)}^{2}\bigr]
  \le m^{-2}L^{2}+2\,m^{-2}L^{2}\sum_{k\neq0}w_2(k)=C(L,m),
\]
the last sum being the one in \eqref{eq:trace-constant} because $w_2(0)=0$.
The series on the left converges because its majorant does, which is (iii).
\end{proof}

\begin{theorem}[Uniform $\HL{2}$ moment]\label{thm:uniform-moment}
Let $L>0$, $m>0$ and let $C(L,m)$ be the constant
\eqref{eq:trace-constant}. For every $N$, $a$ with $aN=L$, with $\mu_N$ the
law \eqref{eq:finite-source-law} of the finite-source lift,
\[
  \E_{\mu_N}\norm{\varphi_N}^2_{\HL{2}}\le C(L,m),
  \qquad\text{hence}\qquad
  \sup_{N}\E_{\mu_N}\norm{\varphi_N}^2_{\HL{2}}\le C(L,m).
\]
The constant is independent of $N$: it is a function of $L$ and $m$ alone.
\end{theorem}

\begin{proof}
By \eqref{eq:energy} the left side is the expectation of a countable sum of
nonnegative continuous functions of the carrier point, a probe being the
stored coordinate divided by the fixed positive number $\sqrt{w_3(k)}$
(\cref{eq:lift}). Tonelli's theorem interchanges sum and expectation, and
\eqref{eq:trace-bound} bounds the result by $C(L,m)$. For the uniformity,
$aN=L$ holds at every index by \cref{hyp:volume} and
\eqref{eq:trace-constant} contains neither $N$ nor $a$ nor $n$, so the
bound is the same at every index and the supremum leaves it unchanged.
\end{proof}

\subsection{The zero mode}

\begin{lemma}[The zero mode]\label{lem:zero-mode}
Let $Z:=\{\varphi:\widehat\varphi(k)=0 \text{ for every } k\neq0\}$ be the
zero-mode line and $p(\varphi):=\sum_k w_2(k)\abs{\widehat\varphi(k)}^2$.
\emph{(i)} For every carrier point $\varphi$,
\[
  \sum_{k\in\Z^2}\widetilde w_2(k)\abs{\widehat\varphi(k)}^2
  =\abs{\widehat\varphi(0)}^{2}
   +\sum_{k\in\Z^2}w_2(k)\abs{\widehat\varphi(k)}^{2}.
\]
\emph{(ii)} $p$ annihilates $Z$ and is not a norm; for every $R\ge0$ the
set $\{p\le R\}$ contains $Z$, whose closure in the carrier is not compact.
\emph{(iii)} The energy \eqref{eq:energy} restricted to $Z$ equals
$\abs{\widehat\varphi(0)}^{2}$.
\end{lemma}

\begin{proof}
(i) is $\widetilde w_2=w_2+\mathbf 1_{\{0\}}$ read termwise, the two cases
being $w_2(0)=0$ and $\widetilde w_2(0)=1$. For (ii): if $\varphi\in Z$
every summand of $p$ vanishes, so $p$ has non-trivial kernel and is only a
seminorm; $Z$ is a closed unbounded line contained in $\{p\le R\}$, and no
closed subset of the carrier containing an unbounded line is compact. For
(iii): on $Z$ only the $k=0$ summand survives, with weight
$\widetilde w_2(0)=1$.
\end{proof}

The lemma forces the two usages. The energy that defines the sets of
\cref{thm:compact-tail} must be \eqref{eq:energy}, with the full weight: by
(ii) a sublevel set of the seminorm $p$ contains the unbounded line $Z$, so
its closure cannot be compact, and tightness built on $p$ would have no
compact set to offer. The zero mode is therefore paid for by the energy, and
the value $\widetilde w_2(0)=1$ is that payment.

The trace in \cref{lem:trace-constant} is the opposite case: the seminorm
weight is used there because the zero mode is settled separately and
exactly. By (i) the full weighted trace splits into the zero-mode moment
plus a $w_2$-weighted remainder, and by step (ii) of
\cref{lem:trace-constant} that moment equals $m^{-2}L^{2}$ rather than
merely being bounded by $2Vm^{-2}$. Writing the constant with $w_2$ and
adding $m^{-2}L^{2}$ back is thus the exact decomposition of the full trace,
with the zero mode counted once and without a factor $2$; applying the
uniform per-mode bound at $k=0$ as well would have lost the sharp value
$V/m^{2}$. So the zero mode is included in the energy, which is what makes
the sets of \cref{thm:compact-tail} compact, and excluded from the weighted
trace, whose first term accounts for it.

\subsection{Compact tails and tightness}

\begin{theorem}[Compact tails]\label{thm:compact-tail}
Let $S\subseteq\HL{2}$ with $\sup_{\varphi\in S}\norm{\varphi}_{\HL{2}}
\le B<\infty$. \emph{(i)} For every $\delta>0$ there is a finite set
$K=K(\delta)\subset\Z^2$ with, termwise,
\[
  \sum_{k\notin K}w_3(k)\abs{\widehat\varphi(k)}^{2}
  \le\delta\sum_{k\notin K}\widetilde w_2(k)\abs{\widehat\varphi(k)}^{2}
  \le\delta B^{2}
\]
for every $\varphi\in S$. \emph{(ii)} The closure of $S$ in the carrier is
compact. In particular this holds for every sublevel set
$\{\varphi:\norm{\varphi}_{\HL{2}}\le R\}$ with $R\ge0$.
\end{theorem}

\begin{proof}
\emph{Step 1: the weight comparison.} For $k\ne0$ the ratio
$w_3(k)/\widetilde w_2(k)=\norm{k}_\infty^{-2}$ tends to $0$ along the mode
lattice, and that is \cref{prop:completion}(ii). For $\delta>0$ put
$K:=\{k:\norm{k}_\infty\le\lceil\delta^{-1/2}\rceil\}$, a box of finitely
many integer vectors containing $0$; for $k\notin K$ the ratio is at most
$\delta$. So the first inequality of (i) is termwise and the second is
\eqref{eq:energy}: a tail bound of
$\delta\norm{\varphi}^2_{\HL{2}}\le\delta B^{2}$ that mentions neither
$\varphi$ nor $N$, one finite mode set serving all of $S$. The quantitative
form of the same estimate is \cref{prop:completion}(iii).

\emph{Step 2: a diagonal extraction.} Let $(\varphi_j)\subset S$. For each
$k$, $\widetilde w_2(k)\abs{\widehat\varphi_j(k)}^2\le B^{2}$ with
$\widetilde w_2(k)>0$, so every coordinate sequence is bounded. Exhaust
$\Z^2$ by increasing finite boxes $K_1\subset K_2\subset\cdots$. On $K_1$
the coordinates lie in a compact subset of $\C^{K_1}$, so a subsequence
converges there; a subsequence of that one converges on $K_2$; and so on.
The diagonal subsequence converges at every $k$: the diagonal argument. Let
$c=(c_k)$ be the coordinate limit.

\emph{Step 3: convergence in the carrier.} Fix $\delta>0$ and $K=K(\delta)$.
Along the diagonal subsequence,
\[
  \norm{\varphi_j-\varphi_{j'}}^2_{\mathcal C}
  =\sum_{k\in K}w_3(k)\abs{\widehat\varphi_j(k)-\widehat\varphi_{j'}(k)}^{2}
   +\sum_{k\notin K}w_3(k)\abs{\widehat\varphi_j(k)-\widehat\varphi_{j'}(k)}^{2},
\]
where the first sum tends to $0$, being a finite sum of convergent
coordinates, and the second is at most $2\delta B^{2}+2\delta B^{2}
=4\delta B^{2}$ by step 1 and $\abs{u-v}^{2}\le2\abs u^{2}+2\abs v^{2}$. As
$\delta$ is arbitrary the subsequence is Cauchy, and the carrier is complete
(\cref{prop:completion}(i)), so it converges to a point of the closure of
$S$. Every sequence in $S$ thus has a subsequence converging in the carrier,
so $S$ is totally bounded, and the closure of a totally bounded set in a
complete metric space is compact: that is (ii).
\end{proof}

\begin{theorem}[Tightness]\label{thm:tightness}
Let $L>0$, $m>0$, and let $\mu_N$ be the law \eqref{eq:finite-source-law}
at $N=n+1$, $a=L/N$. The family $(\mu_N)$ is tight on the carrier: for
every $\epsilon>0$ there is a compact set $K_\epsilon$ in the carrier with
$\mu_N(K_\epsilon)\ge1-\epsilon$ for every $N$.
\end{theorem}

\begin{proof}
Let $C:=C(L,m)$ be \eqref{eq:trace-constant} and put $T:=C/\epsilon+1$,
$K_\epsilon:=\overline{\{\varphi:\norm{\varphi}^2_{\HL{2}}\le T\}}$, the
closure being taken in the carrier; $K_\epsilon$ is compact by
\cref{thm:compact-tail}(ii), applied with $B=\sqrt{T}$. Its complement lies
where the energy exceeds $T$, so Markov's inequality applied to the
extended-valued energy, which is measurable as a countable sum of continuous
functions, with $\E\norm{\varphi_N}^2_{\HL{2}}\le C$ by
\cref{thm:uniform-moment}, gives
\[
  \mu_N(K_\epsilon^{\mathrm c})
  \le\frac{1}{T}\,\E_{\mu_N}\norm{\varphi_N}^2_{\HL{2}}
  \le\frac{C}{T}<\epsilon ,
\]
the last inequality because $T>C/\epsilon$. The set $K_\epsilon$ therefore
satisfies $\mu_N(K_\epsilon)\ge1-\epsilon$, it depends on $\epsilon$, $L$
and $m$ only, and the $+1$ in the cutoff is used exactly here: it makes the
Markov tail strictly smaller than $\epsilon$.
\end{proof}

\begin{external}[Prokhorov]\label{ext:prokhorov}
Let $X$ be a Polish space and $(\mu_N)$ a sequence of Borel probability
measures on $X$. If $(\mu_N)$ is tight, then there are a strictly increasing
sequence $N_1<N_2<\cdots$ and a Borel probability measure $\mu$ on $X$ with
$\int h\,\dd\mu_{N_j}\to\int h\,\dd\mu$ for every bounded continuous
$h:X\to\R$ \cite{prokhorov1956}.
\end{external}

\Cref{ext:prokhorov} produces a subsequence and nothing about the original
sequence: the whole-sequence statement is \cref{thm:main}, reached in
\cref{sec:main} by a separate step and not as a consequence of this
theorem. Its Polish hypothesis is met, the carrier being complete
(\cref{prop:completion}(i)) and separable (\cref{thm:separable}).

One honest limitation closes the section. The only input to
\cref{thm:tightness} is Markov's inequality applied to the second moment of
\cref{thm:uniform-moment}: no exponential moment, no layer-cake
decomposition and no quantitative modulus enters, so the tail bound is
$C/T$ and no better. Consequently this paper asserts no rate of weak
convergence, no exponential tail bound and no modulus of continuity in the
limit, and a reader who wants one must supply an exponential moment for the
family $\mu_N$, which is not proved here. \Cref{sec:scope} records this
among the other boundaries of the result.
